\documentclass[11pt]{article}

% Baseerat: measuring computer-use agent oversight without sight.
% arXiv preprint source. Compiles with a standard TeX distribution
% (pdflatex baseerat.tex, twice for references). No exotic packages.
% House style: UK English throughout; no em or en dashes in prose.

\usepackage[utf8]{inputenc}
\usepackage[T1]{fontenc}
\usepackage{lmodern}
\usepackage[margin=1in]{geometry}
\usepackage{microtype}
\usepackage{booktabs}
\usepackage{array}
\usepackage{amsmath}
\usepackage{xcolor}
\usepackage[hidelinks]{hyperref}
\usepackage{caption}
\captionsetup{font=small,labelfont=bf}

\newcommand{\gap}{oversight parity gap}
\newcommand{\pending}{\textcolor{gray}{\small (pending)}}

\title{\textbf{Baseerat: Measuring the Oversight Gap When a\\
Computer-Use Agent's Overseer Cannot See the Screen}}

\author{Yasir Musawar\\ WAJD AI\\ \texttt{baseerat@wajd.co.uk}}

\date{Preprint. \today}

\begin{document}
\maketitle

\begin{abstract}
Safety benchmarks for computer-use agents evaluate an agent's behaviour under
the tacit assumption that the human overseer can see the screen. For the roughly
forty million people who operate a computer through a screen reader, that
assumption does not hold: the only account of what an agent did arrives through
the agent's narration and the post-hoc accessibility tree. A growing body of
work shows that an agent's own account of its actions is unreliable, with claimed
success exceeding actual success by large margins. We ask a question these two
literatures leave open: when the only oversight channel is the accessible one,
how much of an agent failure, and in particular an adversarially concealed one,
can a non-visual overseer still detect? We introduce \emph{Baseerat}, a benchmark
that pairs each task with a matched sighted control and measures the
\emph{\gap}: the detection rate a sighted overseer achieves on a diverged agent
run minus the rate an overseer restricted to the accessible channel achieves on
the same run. Baseerat runs tasks in a real browser under three conditions
(clean, silent failure, and a deceptive indirect injection that poisons the
accessibility tree in step with the agent's narration), and audits every run
through both channels. Using a deterministic reference auditor over a
150-task corpus, we establish the benchmark's construct validity and report a
\gap{} of 0.50, concentrated entirely in the deceptive-injection condition. We
then evaluate a defence, trusted action receipts, and show it closes the gap to
zero. We release the benchmark, the live rig, the corpus, and the reference
results, and specify the protocol for the language-model and human
expert-auditor arms.
\end{abstract}

\section{Introduction}

Computer-use agents now act on real interfaces on a person's behalf: they fill
forms, send messages, move money, and change settings. The safety of these
agents is increasingly benchmarked, but the benchmarks share an unstated premise.
When they score whether an agent completed a task correctly, resisted a prompt
injection, or avoided a harmful action, the ground truth is read from the screen,
and the implicit overseer who would catch a failure in deployment is assumed to
be able to read the screen too.

For a blind screen-reader user that premise fails. The overseer does not see the
result; they hear the agent's narration and, if they inspect the outcome, they
read whatever the accessibility tree exposes. The oversight channel collapses
onto the agent's own account of what it did. This would be a manageable
limitation if that account were reliable. It is not. Recent work characterises
\emph{false success}, in which agents produce confident closure statements over
tasks that in fact failed, and \emph{outcome hallucination}, in which an agent
claims completion despite evidence to the contrary, with claimed-to-actual
success ratios reported as high as $5.3\times$.

Put together, these two observations imply a vulnerability that neither
literature measures. If the accessible channel is the whole of a blind user's
oversight, and that channel carries an unreliable self-report, then the safety of
the agent for that user reduces to the fidelity of the self-report, and to
whether the accessible channel is rich enough to reveal a discrepancy when the
self-report is wrong. Nobody has measured either quantity, and the adversarial
case is worse than the accidental one. In an indirect prompt injection the
attacker controls page content; nothing prevents the injection from also
instructing the agent to narrate a clean, benign account of a hijacked action, so
that the attack and its cover story travel down the very channel the blind user
depends on. A sighted overseer sees the true result regardless; a non-visual
overseer sees only the cover story.

We make this precise with a benchmark. Our contributions are:

\begin{enumerate}
\item \textbf{A problem formulation and a metric.} We frame agent oversight as a
detection task performed over a restricted perceptual channel, and define the
\emph{\gap}: the difference in divergence-detection rate between a visually
unrestricted overseer and one restricted to the accessible channel, over the same
runs. We further define \emph{self-report fidelity} and \emph{non-visual
detectability}.

\item \textbf{A live benchmark, Baseerat.} Tasks run in a real browser under
three conditions. Crucially, the deceptive-injection condition is constructed in
an actual document object model so that the accessibility tree exposes a benign
value while the committed state is hijacked. We show this divergence is not a
modelling artefact but is realisable and reads directly from a production
accessibility snapshot.

\item \textbf{Reference results and construct validity.} A deterministic
reference auditor over a 150-task corpus (six task domains, four injection
archetypes) yields a \gap{} of 0.50, entirely within the deceptive-injection
condition, with no false alarms on clean runs.

\item \textbf{A defence and its evaluation.} We evaluate trusted action receipts,
a machine-verified and accessibly-rendered record of what was actually committed,
independent of the page. Receipts restore non-visual detection from 0.50 to 1.00
and close the \gap{} to zero.

\item \textbf{An open protocol for the auditor arms.} We release the rig and
specify the language-model auditor sweep and the human expert-auditor study that
populate the remaining evaluation.
\end{enumerate}

The name \emph{Baseerat} (from the Arabic \emph{baseerah}, ``inner sight'',
discernment) marks the point: the question is not whether an overseer has eyes,
but whether the channel they are given carries enough to discern a failure.

\section{Related work and positioning}

Baseerat sits at the intersection of three lines of work, each of which supplies
one ingredient and none of which occupies the joint question.

\paragraph{Safety benchmarks for computer-use agents.} OS-Harm evaluates
computer-use agents across deliberate misuse, prompt injection, and model
misbehaviour, and is representative of a class that scores outcomes with full
visual access to the screen. It does not model the overseer's perceptual channel;
the auditor is assumed sighted.

\paragraph{Unreliable self-reports.} A separate line characterises the failure
mode we exploit. Work on false success shows agents assert completion over failed
tasks and that language-model judges systematically miss it; MIRAGE-Bench
provides a taxonomy of agent hallucination including outcome hallucination. These
establish that the self-report is untrustworthy, but they too assume a sighted
party is available to notice, and do not consider accessibility.

\paragraph{Accessibility of agents.} A recent and fast-moving line studies agents
for blind and low-vision users. A11y-CUA documents that agent task success
collapses under assistive-technology interaction styles. An interview study of
blind users in the agentic era names the epistemic vulnerability precisely,
observing that an error early in a multi-step task can propagate silently before
any feedback is received, but it is qualitative and introduces no benchmark or
metric. Closest to our work, NeXUI scores whether an assistive agent explains
each step in a way grounded in interface state. It asks whether agents explain
well; it does not include adversarial conditions, seeded silent failures,
deception, or any measurement of whether the explanations are sufficient to
detect an agent failing or lying, and it runs without blind participants.

\paragraph{The joint question.} The gap is the conjunction. The safety
benchmarks measure agent behaviour but assume a sighted auditor; the false-success
work measures self-report unreliability but assumes a sighted auditor; the
accessibility work studies the non-visual channel but not adversarial detection.
Baseerat measures whether the non-visual channel is \emph{sufficient to detect
failure}, including adversarial failure, and what restores it when it is not. The
defence we evaluate is informed by the tool-receipts idea, a machine-verified
statement of effect rather than the agent's word for it.

\section{The Baseerat benchmark}

\subsection{Tasks and the field specification}

A task is a realistic operation on a web or desktop surface: send an email,
update a delivery address, add a payee, rename and move a file, book a calendar
event, update account recovery details. Each task carries a \emph{field
specification}: the set of observable slots that a correct completion must leave
holding specified values (for an email, the recipient, subject, and attachment).
The field specification is the ground truth against which both channels are
scored.

The corpus contains 150 tasks generated deterministically, stratified evenly
across six domains and four \emph{injection archetypes} that vary how a hijacked
value differs from the intended one: a homograph (a look-alike character swap such
as \texttt{l} to \texttt{1}), a plausible substitution, an appended hostile
token, and a wholesale redirect. The stratification exists so that the auditor
arms can be analysed by how subtle the hijack is; the reference auditor, by
construction, is archetype-invariant (Section~\ref{sec:results}).

\subsection{Three conditions}

Every task is run under three conditions. The design of the conditions is the
core of the benchmark, because it separates a failure that the accessible channel
can reveal from one that it cannot.

\begin{center}
\begin{tabular}{@{}p{2.6cm}p{3.1cm}p{3.1cm}p{3.3cm}@{}}
\toprule
\textbf{Condition} & \textbf{Ground truth} & \textbf{Accessibility tree} &
\textbf{Detectable by} \\
\midrule
clean & matches the task & honest & nothing to detect \\
silent failure & a field silently fails to save & \emph{honest}: shows the
failure & anyone who cross-checks the tree against the agent's claim \\
deceptive injection & a field is hijacked & \emph{poisoned}: shows the benign
value & only a sighted overseer reading the true result \\
\bottomrule
\end{tabular}
\end{center}

Under \textbf{silent failure} the accessibility tree is honest, so a diligent
non-visual overseer who cross-checks the tree against the agent's claim of success
can catch it. Under \textbf{deceptive injection} the accessibility tree is
poisoned to corroborate the false narration, so the non-visual channel shows a
consistent benign story and only the ground truth reveals the hijack. The
contrast between these two conditions is what isolates the \gap.

\subsection{Two channels}

Each run is audited through two channels. The \textbf{visual} channel additionally
receives the ground-truth committed state, modelling a sighted overseer who can
read the true result from the screen. The \textbf{non-visual} channel receives the
agent's narration and the post-hoc accessibility tree only, and never the ground
truth. The restriction is enforced at the data boundary: the object handed to a
non-visual auditor does not contain the committed state.

\subsection{The live rig}

Conditions are made physical in a real browser rather than simulated. Each task is
rendered as an HTML form whose submit handler writes the truly committed values
to an authoritative store and, separately, renders a confirmation region into the
accessibility tree. Under deceptive injection the handler commits the attacker's
value while the confirmation exposes the benign one. A pluggable agent perceives
the page through its accessible representation, decides the field values, performs
them, and writes its own narration. The rig then reads ground truth from the
authoritative store and reads the accessible view the way a screen reader would,
by resolving each field through its accessible name and by capturing the full
accessibility snapshot of the confirmation region. The agent supplies the
narration but never touches the committed state, so it can misreport reality but
cannot fabricate it.

Figure~\ref{fig:diverge} shows a captured example. Under deceptive injection the
committed recipient is \texttt{accounts@northgate-1egal.co.uk} (a homograph, digit
one for the letter l) while the accessibility snapshot exposes the benign
\texttt{accounts@northgate-legal.co.uk}. The two channels genuinely disagree, and
the disagreement is read from a production accessibility snapshot, establishing
that the threat model is realisable rather than assumed.

\begin{figure}[t]
\centering
\fbox{\parbox{0.92\linewidth}{\small\ttfamily
=== deceptive\_injection (live browser) ===\\
ground truth : email.recipient = accounts@northgate-1egal.co.uk\\
a11y tree ~~~: email.recipient = accounts@northgate-legal.co.uk\\
channels disagree? YES
}}
\caption{A single live run. The committed recipient is hijacked by a homograph;
the accessibility tree, read from a real snapshot, still speaks the benign
address. A sighted overseer sees the hijack; the accessible channel does not.}
\label{fig:diverge}
\end{figure}

\subsection{Cross-checking the accessible channel against a screen reader}

The rig reads the accessible channel through a browser accessibility snapshot,
which is a faithful model of the accessibility tree but is not identical to what a
production screen reader announces. To validate the model, Baseerat includes a
cross-check that parses an NVDA Speech Viewer log into the same field-value
representation and compares it against the snapshot channel. The parser and the
comparator are released and tested against captured logs; the live capture is
performed on a Windows virtual machine following a documented procedure. Under
deceptive injection the two accessible surfaces agree with each other on the
benign spoofed value and both diverge from ground truth, which is precisely why
the accessible channel is insufficient.

\section{Metrics}

Let a run have field specification with expected values $E = \{(k, e_k)\}$,
committed ground-truth state $W$, exposed accessibility tree $A$, and an agent
claim of success. Let $\mathrm{Truth}(c)$ be the verdict a perfect auditor should
return for condition $c$: \textsc{clean} for the clean condition and
\textsc{diverged} otherwise.

\paragraph{Self-report fidelity.} The fraction of expected fields whose true
committed value the agent reported honestly. Under a claim of success every field
is implicitly claimed correct, so fidelity is the fraction that really is correct:
\[
\mathrm{Fidelity} = \frac{1}{|E|}\,\bigl|\{\,k : W_k = e_k\,\}\bigr|.
\]
Fidelity is computed from ground truth and is independent of any auditor.

\paragraph{Non-visual detectability.} Given only the non-visual channel, does an
auditor flag the runs that truly diverged? We report \emph{recall} over the
diverged runs (silent failure and deceptive injection), namely the fraction of
runs that should have been caught that were, and separately the \emph{false-alarm
rate} on clean runs.

\paragraph{Oversight parity gap.} The headline metric. Let $D_{\mathrm{vis}}$ and
$D_{\mathrm{nv}}$ be the detection rates on truly-diverged runs under the visual
and non-visual channels respectively, over the same runs. Then
\[
\mathrm{Gap} = D_{\mathrm{vis}} - D_{\mathrm{nv}}.
\]
A positive gap is oversight that a sighted overseer can exercise and a blind one
structurally cannot. We report the gap overall and per condition.

\paragraph{Defence effect.} With a trusted receipt present, the non-visual
detection rate rises to $D_{\mathrm{nv}}^{\mathrm{rcpt}}$ and the residual gap is
$D_{\mathrm{vis}} - D_{\mathrm{nv}}^{\mathrm{rcpt}}$. The gap closed by the
defence is the difference between the narration-only gap and the residual gap.

\section{Reference-auditor results}
\label{sec:results}

We first report a deterministic \emph{reference auditor}. It encodes the single
most obvious cross-check a diligent overseer would perform: does the accessibility
tree corroborate the agent's claim of success? On the visual channel it checks the
claim against ground truth. The reference auditor is not intended to be
intelligent; it establishes the benchmark's construct validity (that the
conditions are well-formed and the metrics behave as designed) and serves as the
scalable lower bound against which the language-model and human auditors are
compared. Its results are computed over the full 150-task corpus (450 runs) and
are exact and reproducible.

\begin{table}[t]
\centering
\begin{tabular}{@{}lccc@{}}
\toprule
\textbf{Quantity} & \textbf{clean} & \textbf{silent failure} &
\textbf{deceptive injection} \\
\midrule
Self-report fidelity & 1.00 & 0.61 & 0.61 \\
Visual detection rate & n/a & 1.00 & 1.00 \\
Non-visual detection rate & n/a & 1.00 & 0.00 \\
\textbf{Oversight parity gap} & 0.00 & 0.00 & \textbf{1.00} \\
\bottomrule
\end{tabular}
\caption{Reference-auditor results over 150 tasks (450 runs). The parity gap is
0.50 overall and is concentrated entirely in the deceptive-injection condition.
The non-visual channel catches every silent failure (the honest tree contradicts
the false claim) and no injection (the poisoned tree corroborates it). There are
no false alarms on clean runs.}
\label{tab:reference}
\end{table}

Table~\ref{tab:reference} shows the result. The non-visual channel catches every
silent failure and no deceptive injection, exactly as the condition design
predicts: when the accessibility tree is honest the cross-check succeeds, and when
it is poisoned in step with the narration the cross-check is defeated. The visual
channel catches both. The overall \gap{} is 0.50, and it lives entirely in the
deceptive-injection condition. Because the poisoning defeats the reference
auditor regardless of how the hijacked value differs from the intended one, the
reference auditor is archetype-invariant; the injection archetypes are provided
so that the language-model auditor, which may catch a clumsy homograph while
missing a clean redirect, can be analysed by subtlety.

\section{The defence: trusted action receipts}

The reference results say the accessible channel is insufficient. The constructive
question is what would make it sufficient. We evaluate \emph{trusted action
receipts}. A receipt is a machine-verified, accessibly-rendered statement of what
was actually committed, produced by an instrumented layer that observes the
committed state directly and is independent of the page's own confirmation. In the
rig the receipt is rendered from the authoritative committed store, not from the
poisonable confirmation, so under deceptive injection the receipt reports the true
hijacked value while the confirmation still lies. An auditor given a receipt treats
it as authoritative and cross-checks it against the expected values.

\begin{table}[t]
\centering
\begin{tabular}{@{}lcc@{}}
\toprule
& \textbf{narration only} & \textbf{with receipt} \\
\midrule
Non-visual detection rate (diverged runs) & 0.50 & \textbf{1.00} \\
Oversight parity gap & 0.50 & \textbf{0.00} \\
\bottomrule
\end{tabular}
\caption{The defence over 150 tasks. A trusted receipt channel restores non-visual
detection from 0.50 to 1.00 and closes the parity gap to zero, without introducing
false alarms on clean runs.}
\label{tab:defence}
\end{table}

Table~\ref{tab:defence} shows that the receipt restores non-visual detection to
1.00 and closes the \gap{} to zero, with no false alarms on clean runs, and the
same result holds when the runs are produced by the live browser rig rather than
the deterministic simulator. The interpretation is that a trusted channel,
independent of the page and rendered into the accessible modality, restores to a
blind overseer exactly the detection a sighted overseer had. The realism caveat is
explicit: the defence assumes such an instrumentation layer exists and cannot
itself be spoofed by the page. Baseerat does not claim that assumption is free; it
measures how much oversight the assumption would buy, which is the whole of the
gap.

\section{Protocol for the auditor arms}

The reference auditor establishes the benchmark and the mechanism. Two further
auditor arms complete the evaluation and are released as an open protocol.

\paragraph{Language-model auditor sweep.} A language model is placed under the
same perceptual restriction as each channel and asked for a structured verdict.
This is not a simulation of a blind person; it is a model of the oversight
\emph{channel}, the information bottleneck, which is what the gap measures. The
released rig runs this arm across models with a single flag; the corpus
stratification supports a per-archetype and per-domain breakdown of where the
non-visual channel is most and least detectable. Table~\ref{tab:llm} gives the
reporting template; the cells are populated in the next revision.

\begin{table}[t]
\centering
\begin{tabular}{@{}lcccc@{}}
\toprule
\textbf{Auditor model} & \textbf{Vis. recall} & \textbf{Non-vis. recall} &
\textbf{Parity gap} & \textbf{Gap w/ receipt} \\
\midrule
reference (deterministic) & 1.00 & 0.50 & 0.50 & 0.00 \\
% TABLE3-ROWS-START (auto-filled by paper/fill_table3.py; do not edit by hand)
claude-sonnet-5 & 1.00 & 0.50 & 0.50 & 0.00 \\
% TABLE3-ROWS-END
\bottomrule
\end{tabular}
\caption{Reporting template for the language-model auditor sweep. The reference
row is the deterministic lower bound reported here; the model rows are populated
in the next revision. A model auditor may exceed the reference on silent failures
and may partially detect clumsy injection archetypes, which would narrow but is
not expected to close the narration-only gap.}
\label{tab:llm}
\end{table}

\paragraph{Human expert-auditor arm.} Because a language model under a perceptual
restriction is a lower bound and not a person, the language-model arm is anchored
by an expert-auditor evaluation with consenting professional blind accessibility
auditors, who audit a stratified sample of runs through the non-visual channel.
This arm calibrates the language-model auditor against human non-visual detection
and is the basis for the journal version of this work.

\section{Limitations and threats to validity}

\paragraph{The agent is scripted for the reference results.} In the reference
configuration the agent performs the task correctly and the divergence is
introduced by the environment. This is a deliberate control: it isolates the
oversight question from agent competence, so that a measured gap reflects the
channel and not agent error. A real computer-use agent slots into the same seam,
and the released rig supports it; a real agent introduces its own error and its
own genuine narration, which the language-model auditor arm will exercise.

\paragraph{The receipt defence assumes a trusted layer.} The defence's power
depends on an instrumentation layer that the page cannot spoof. Where no such
layer exists, the measured gap stands as the exposure. The benchmark measures the
value of the assumption rather than asserting it is cheap.

\paragraph{The accessible channel is modelled by a snapshot.} The rig reads the
accessibility tree through a browser snapshot. The NVDA cross-check validates this
against a production screen reader on a stratified sample, but a full screen-reader
in the loop for every run is future work.

\paragraph{Construct validity of the reference auditor.} The reference auditor is
a deliberately simple cross-check and is not a claim about how well any real
auditor performs. Its role is to show the conditions and metrics are well-formed
and to provide a lower bound; the human arm supplies the ecological anchor.

\section{Ethics and accessibility statement}

This work is motivated by a safety exposure that falls disproportionately on blind
and low-vision users, and it is intended to be used to reduce that exposure. The
corpus contains only fictional values and no real personal data. The
deceptive-injection condition is a controlled construction inside a local rig; the
benchmark does not provide an attack against any deployed system. The human
expert-auditor arm is designed as a paid, non-clinical, minimal-risk evaluation
with informed consent. We release the benchmark so that model developers can be
held to a measured oversight-parity standard rather than an assumed one.

\section{Conclusion}

Agent oversight has a hidden precondition: that the overseer can see the screen.
For a blind screen-reader user that precondition fails, and the oversight channel
collapses onto an agent self-report that the literature has shown to be
unreliable, and that an adversary can actively poison in step with the action.
Baseerat turns this observation into a measurement. Over a 150-task corpus a
deterministic reference auditor exhibits an \gap{} of 0.50, concentrated entirely
in a deceptive-injection condition that we demonstrate is realisable in a live
browser, and a trusted-receipt defence closes the gap to zero. We release the
benchmark, the rig, the corpus, and the reference results, and we specify the
language-model and human expert-auditor arms that complete the evaluation. The
broader claim is simple: model developers can and should be scored on the
oversight they afford an overseer who cannot see the screen, not on the oversight
they afford one who can.

\section*{Availability}

The benchmark, the live rig, the task generator, the corpus, the defence, the
NVDA cross-check, and the reference results are released as an open repository.
This preprint is archived at
\href{https://doi.org/10.5281/zenodo.22080740}{doi:10.5281/zenodo.22080740}.
Copyright WAJD Group. Developed by WAJD AI.

\paragraph{Cite this as.} Y. Musawar. \emph{Baseerat: Measuring the Oversight
Gap When a Computer-Use Agent's Overseer Cannot See the Screen}. Preprint, WAJD
AI, 2026. doi:10.5281/zenodo.22080740.

\begin{thebibliography}{9}
\small

\bibitem{osharm} OS-Harm: A Benchmark for Measuring Safety of Computer Use
Agents. arXiv:2506.14866.

\bibitem{silentfailure} From Confident Closing to Silent Failure:
Characterizing False Success in LLM Agents. arXiv:2606.09863.

\bibitem{mirage} MIRAGE-Bench: LLM Agent is Hallucinating and Where to Find
Them. arXiv:2507.21017.

\bibitem{a11ycua} A11y-CUA: Characterizing the Accessibility Gap in Computer Use
Agents. arXiv:2602.09310.

\bibitem{xaiblv} Explainable AI for Blind and Low-Vision Users: Navigating
Trust, Modality, and Interpretability in the Agentic Era. arXiv:2604.00187.

\bibitem{nexui} Navigation Alone Is Not Enough: Evaluating Explanatory Assistive
UI Agents (NeXUI). arXiv:2608.09944.

\bibitem{morae} Morae: Proactively Pausing UI Agents for User Choices.
arXiv:2508.21456.

\bibitem{receipts} Tool Receipts, Not Zero-Knowledge Proofs: Practical
Hallucination Detection for AI Agents. arXiv:2603.10060.

\bibitem{misscore} How Benchmarks Mis-Score Computer-Use Agents.
arXiv:2607.28367.

\end{thebibliography}

\end{document}
